(1) For $v \in V_\C$ define the bilinear form $G_v$ on $\CSm$ by
\[
G_v(s_1, s_2) \;:=\; B\bigl(s_1,\, c(v)\,s_2\bigr), \qquad s_1, s_2 \in \CSm,
\]
which is well defined since $c(v)s_2 \in \CS$ and $B$ pairs $\CSm$ with $\CS$ (Lemma \claimref{correspondence.chevalley-pairing}(1)). $G_v$ is antisymmetric: $G_v(s_2, s_1) = B(s_2, c(v)s_1) = B(c(v)s_2, s_1) = -B(s_1, c(v)s_2) = -G_v(s_1,s_2)$ by clauses (3) and (2) of the lemma. The perfect pairing $B$ identifies $\CSm^{\,*}$ with $\CS$, hence $\Lambda^2(\CSm^{\,*})$ with $\Lambda^2\CS$; let $\Psi(v) \in \Lambda^2\CS$ be the image of $G_v$. The map $\Psi: V_\C \to \Lambda^2\CS$ is linear, and it is injective: if $\Psi(v) = 0$ then $c(v)$ vanishes on $\CSm$ ($B$ is perfect). Choose a non-null $w \perp v$ (the orthogonal complement of $v$ has dimension at least $5$ and a totally null subspace has dimension at most $3$). For $s \in \CSm$, $c(v)c(w)s = 2\ip{v}{w}s - c(w)c(v)s = 0$ (\claimref{std-algebra.vector-squares}), and $c(w): \CSm \to \CS$ is bijective because $c(w)^2 = Q(w) \neq 0$ (\claimref{std-algebra.inverse-of-vector}); so $c(v)$ vanishes on $\CS$ as well, $c(v) = 0$, and $v = 0$ by injectivity of $c$ (Theorem \claimref{correspondence.fock-model-complexified-ascended-algebra}). Since $\dim_\C \Lambda^2\CS = \binom42 = 6 = \dim_\C V_\C$, $\Psi$ is an isomorphism. Equivariance: $\mathrm{Spin}$ acts on $c(v)$ by conjugation, giving $c(gv)$, and on $G_v$ through its module actions; $B$ is invariant (Lemma \claimref{correspondence.chevalley-pairing}(4)), so $\Psi(gv) = g\,\Psi(v)$ with the induced action on $\Lambda^2\CS$. The form: $(v, v') \mapsto \varepsilon_0(\Psi(v)\wedge\Psi(v'))$, with $\varepsilon_0$ any generator of $(\Lambda^4\CS)^*$, is a symmetric $\mathrm{Spin}$-invariant bilinear form on $V_\C$, hence proportional to $\ip{\cdot}{\cdot}$ (the space of invariant symmetric forms on the irreducible $V_\C$ is one-dimensional: an invariant form is $\ip{T\cdot}{\cdot}$ for a $T$ commuting with the action, hence a scalar by Schur; and $V_\C$ is irreducible under $\mathfrak{so}(V_\C)$, since for an invariant $U \ni v \neq 0$ with $v_j := \ip{e_j}{v} \neq 0$ in an orthonormal basis, $(e_k\wedge e_i)(e_i\wedge e_j)v = Q(e_i)\,v_j\,e_k$ for distinct $i, j, k$ and $(e_j\wedge e_k)e_k = Q(e_k)\,e_j$, nonzero multiples in either case, so $U = V_\C$). The constant is not zero and is computed in the standard bases: with $\CS$-basis $t_0 = 1$, $t_1 = w_2\wedge w_3$, $t_2 = w_3 \wedge w_1$, $t_3 = w_1\wedge w_2$, $\CSm$-basis $(w_1, w_2, w_3, w_{123})$, and $\varepsilon_0(t_0\wedge t_1\wedge t_2\wedge t_3) = 1$, a direct computation (printed in the script, and in Example \ref{ex:psi} below) gives $\Psi(w_1) = -\sqrt2\; t_2\wedge t_3$ and $\Psi(w_1') = \sqrt2\; t_0 \wedge t_1$, whence $\varepsilon_0(\Psi(w_1)\wedge\Psi(w_1')) = -2 = -2\ip{w_1}{w_1'}$. Setting $\varepsilon := -\tfrac12\,\varepsilon_0$ gives $\varepsilon(\Psi(v)\wedge\Psi(v')) = \ip{v}{v'}$ exactly, which the script confirms on all twenty-one basis pairs.
(2) A bivector $P \in \Lambda^2\CS$ on a four-dimensional space is decomposable if and only if $P \wedge P = 0$: decomposable implies $s\wedge t\wedge s\wedge t = 0$; conversely, an antisymmetric bilinear form on $\CS^*$ has a basis in which it reads $t_0\wedge t_1$ or $t_0\wedge t_1 + t_2\wedge t_3$ (pick $t_0, t_1$ with $P(t_0,t_1) = 1$ and repeat on their $P$-orthogonal complement), and in the second case $P\wedge P = 2\,t_0\wedge t_1\wedge t_2\wedge t_3 \neq 0$; in coordinates $P\wedge P = 2\,\mathrm{Pf}(P)\, t_0\wedge t_1\wedge t_2 \wedge t_3$ with $\mathrm{Pf}$ the Pfaffian, which the script uses. By (1), $Q(v) = \varepsilon(\Psi(v)\wedge\Psi(v))$, so $v$ is null iff $\Psi(v)\wedge\Psi(v) = 0$ iff $\Psi(v)$ is simple, i.e.\ represents a $2$-plane of $\CS$, i.e.\ a line of $\PCS$.
(3) The fibre quadric (\claimref{correspondence.imported-forced-completion-fibre-quadric}) is the space of null rays of the real slice $V \subset V_\C$; each null ray spans a complex null line, hence by (2) a point of the Klein quadric $\{[P] : P\wedge P = 0\} \subset \mathbb{P}(\Lambda^2\CS)$, i.e.\ a line of $\PCS$. The assignment factors through the null lines of $V$ (a ray and its opposite span the same complex line) and is injective on them. Which lines of $\PCS$ arise from real points is a reality condition, settled in Proposition \claimref{correspondence.reality-conjugation-twistor-form-factor}(5).
